% Blank article -- SkoLab starter (CC0: use it for anything).
\documentclass[11pt]{article}
\usepackage[utf8]{inputenc}
\usepackage[T1]{fontenc}
\usepackage{lmodern}
\usepackage[margin=1in]{geometry}
\usepackage{amsmath,amssymb}
\usepackage{graphicx}
\usepackage{booktabs}
\usepackage[hidelinks]{hyperref}

\title{Your Title Here}
\author{First Author \and Second Author}
\date{\today}

\begin{document}

\maketitle

\begin{abstract}
A short summary of the problem, what you did, and what you found.
\end{abstract}

\section{Introduction}
\label{sec:intro}
Write your introduction here. Cite related work like this~\cite{lamport94}
and refer to sections like this: Section~\ref{sec:method}.

\section{Method}
\label{sec:method}
Equations are numbered automatically:
\begin{equation}
  \label{eq:energy}
  E = mc^2 .
\end{equation}
Equation~\eqref{eq:energy} is referenced by its label.

\begin{table}[h]
  \centering
  \begin{tabular}{lrr}
    \toprule
    Setting & Accuracy & Time (s) \\
    \midrule
    Baseline & 0.81 & 12.4 \\
    Ours     & 0.87 & 10.9 \\
    \bottomrule
  \end{tabular}
  \caption{An example table.}
  \label{tab:results}
\end{table}

\section{Conclusion}
Summarize your contributions.

\begin{thebibliography}{9}
\bibitem{lamport94}
  Leslie Lamport,
  \emph{\LaTeX: A Document Preparation System},
  Addison-Wesley, 2nd edition, 1994.
\end{thebibliography}

\end{document}
